% Generated draft: requires mathematical and source review.
\documentclass[11pt]{article}
\usepackage[margin=25mm]{geometry}
\usepackage{fontspec}
\setmainfont{Latin Modern Roman}
\usepackage{amsmath,amssymb,mathtools,unicode-math}
\setmathfont{Latin Modern Math}
\usepackage{enumitem,needspace,xurl,hyperref}
\setlength{\parindent}{0pt}
\setlength{\parskip}{5pt}
\setlength{\emergencystretch}{3em}
\newcommand{\N}{\mathbb{N}}
\newcommand{\Z}{\mathbb{Z}}
\newcommand{\Q}{\mathbb{Q}}
\newcommand{\R}{\mathbb{R}}
\newcommand{\C}{\mathbb{C}}
\newcommand{\F}{\mathbb{F}}
\newcommand{\A}{\mathbb{A}}
\newcommand{\PP}{\mathbb{P}}
\newcommand{\E}{\mathbb{E}}
\newcommand{\eps}{\varepsilon}
\newcommand{\dd}{\,\mathrm{d}}
\begin{document}
\section*{UnsolvedMath ID: 30002342}
\textbf{Source problem code:} catalogue entry
\subsection{Definitions and mathematical statement}
Let \(X\) be a singular complex algebraic variety, let \(U=X_{\mathrm{reg}}\) denote its smooth locus, and let \(j:U\hookrightarrow X\) be the inclusion. The source imposes no explicit normality, irreducibility, connectedness, or properness hypothesis; none is added here.

For a complex algebraic variety \(T\), let \(\operatorname{FinEt}(T)\) be the category of finite étale morphisms to \(T\), with morphisms over \(T\). Given a finite étale cover \(p:V\to U\), an extension as a finite étale cover of \(X\) consists of a finite étale morphism \(\overline p:Y\to X\) and an isomorphism
\[
Y\times_X U\cong V
\]
over \(U\). Thus restriction defines a functor
\[
j^*:\operatorname{FinEt}(X)\longrightarrow\operatorname{FinEt}(U).
\]
For a geometric base point \(\bar x\) of \(U\), the étale fundamental group of a connected component containing \(\bar x\) is the profinite automorphism group of the geometric-fiber functor on its finite étale covers. Writing \(U_{\bar x}\) and \(X_{\bar x}\) for the respective connected components, inclusion induces a continuous homomorphism
\[
j_*:\pi_1^{\mathrm{\acute et}}(U_{\bar x},\bar x)
\longrightarrow
\pi_1^{\mathrm{\acute et}}(X_{\bar x},\bar x).
\]

\textbf{Problem.} Characterize the essential image of \(j^*\): determine the obstructions to a finite étale cover of \(X_{\mathrm{reg}}\) admitting an extension as a finite étale cover of \(X\). Describe how the singularities of \(X\) affect the comparison homomorphisms \(j_*\), including their kernels and images, and explain their relationship to these extension obstructions. If the smooth locus has several connected components, the extension question concerns the full restriction functor, including compatibility across those components.

This is a formulation of the source's open-ended comparison question, not a proposed theorem. Extension is interpreted here as extension that remains finite étale over all of \(X\); the source does not explicitly specify whether a broader class of extensions was intended. Allowing ramification over the singular locus, including quasi-étale extensions in an appropriate normal setting, is a different extension question. The intended generality of \(X\), the permitted extensions, and the desired form of an obstruction theory require source clarification.
\subsection{Short English statement}
For a singular complex algebraic variety, determine which finite étale covers of its smooth locus extend to finite étale covers of the whole variety, and describe the induced comparison of étale fundamental groups. The source leaves the hypotheses and the intended obstruction theory unspecified. Its supplied status is \emph{partially solved}; this draft does not independently verify that assessment.
\subsection{Sources}
\begin{itemize}
\item Local UnsolvedMath registry, permanent ID \(30002342\), title \emph{Étale Covers of Smooth Loci of Singular Varieties}, exact source problem code \texttt{catalogue entry}, in the \emph{Oberwolfach Reports Open Problems} collection. No individual proposer is supplied; the proposed year is recorded as 2013.
\item Supplied report attribution: \emph{Complex Algebraic Geometry (2014)}, Oberwolfach Reports, DOI \texttt{10.4171/owr/2013/27}, pp.~1563--1627. \url{https://doi.org/10.4171/owr/2013/27}. The supplied background describes these reports as records of mathematical workshops and their open-problem discussions. The report title's year 2014 differs from the recorded proposed year and the year embedded in the DOI; this draft retains all three without resolving the discrepancy.
\item Greb--Kebekus--Peternell, \emph{Etale fundamental groups of klt spaces}. Supplied literature link: \url{https://arxiv.org/abs/1307.5718}.
\item Xu, \emph{Finiteness of algebraic fundamental groups}. Supplied literature link: \url{https://arxiv.org/abs/1210.5564}.
\item Supplied literature assessment, checked on \texttt{2026-08-21}: the registry labels the problem \emph{partially solved} and reports that, for klt spaces, finite quasi-étale covers can be chosen so that the fundamental groups of the cover and its smooth locus agree, and that local algebraic fundamental groups are finite. It further reports that the full direct comparison and obstruction theory requested here is not known in general. This is supplied dated curation derived from the OWR record, not an independently verified literature assessment. The summary does not specify all hypotheses or which notion of fundamental group is intended in its equality claim; the cited claims and time-sensitive status require independent verification.
\item The registry supplies the statement label \texttt{corrected\_verified} and describes a manual semantic, typographic, and LaTeX review of a meaning-preserving cleanup. These are supplied editorial assertions, not independent verification of mathematical correctness or of the report attribution.
\end{itemize}
\end{document}
